
\section{Section 1: Ratios}

\subsection*{Additional: Ratios and Fractions}
\begin{itemize}
    \item \kr{비는 부분과 전체의 관계를 나타내기 때문에 당연하게도 분수로 변환할 수 있다} (e.g., $2:3 = \frac{2}{3}$)
    \item \kr{비도 분수와 같이 약분할 수 있다} (e.g., $8:20 = 4:10 = 2:5$)
    \item If $a:b = c:d$, then $ad = bc$ (e.g., $4:13 = 8:26 \Rightarrow 4 \times 26 = 104 = 8 \times 13 = 104$)
    \item \kr{이를 활용하여 문제를 풀 수 있다}:
    \[
    2:7 = 5:x \Rightarrow 7 \times 5 = 35,\quad 2x = 35 \Rightarrow x = \frac{35}{2}
    \]
\end{itemize}

\section{Section 2: Rates}

\subsection*{Additional: Distance, Speed, Time}
\kr{거속시; 생각보다 많이 나옴}

\[
\begin{array}{c}
\text{거리} = \text{속도} \times \text{시간} \\
\text{속도} = \frac{\text{거리}}{\text{시간}} \\
\text{시간} = \frac{\text{거리}}{\text{속도}}
\end{array}
\]

\subsection*{Metric Prefixes}
\kr{알아두면 좋은 prefix(접두사) 몇가지}:
\begin{itemize}
    \item $10^9$ = giga = G
    \item $10^6$ = mega = M
    \item $10^3$ = kilo = k
    \item $10^{-2}$ = centi = c
    \item $10^{-3}$ = milli = m
    \item $10^{-6}$ = micro = $\mu$
    \item $10^{-9}$ = nano = n
\end{itemize}

\kr{예시: 1m (미터)를 기준으로 1mm = 0.0001m}

\section{Section 3: Unit Rate}

\subsection*{Additional: Complex Fraction}
\kr{Complex fraction(번분수)는 특히 unit rate에서 자주 나옴}

\kr{번분수의 계산은 $a/b \div c/d$ 일 경우 $\frac{ad}{bc}$}

\kr{예시: 서점에서 \$\frac{19}{4}에 5권의 책을 팔 때 1권의 가격은?}

\[
\frac{19}{4} \div 5 = \frac{19}{4} \div \frac{5}{1} = \frac{19 \times 1}{4 \times 5} = \frac{19}{20}
\]

\section{Section 4: Proportions}

\subsection*{Additional: Butterfly Method}
\kr{Cross multiplication(교차 곱셈)의 또 다른 이름}

\[
\frac{x}{5} = \frac{8}{2} \Rightarrow x \times 2 = 5 \times 8 \Rightarrow 2x = 40 \Rightarrow x = 20
\]

\kr{모양이 나비 같다 하여 butterfly method라 이름을 붙였는데 솔직히 억지 같긴함}

\section{Section 5: Percents}

\subsection*{Additional: Decimal and Percentage}
\kr{너무 당연하겠지만 decimal(소수) $\times 100$ = percentage(백분율); vice versa}

\kr{소수나 백분율 관련문제는 대표적으로 가격 할인이나 인상이 있음}

\kr{예시: 정가가 \$100인 옷을 정가의 30\%를 인상해 판매 후 연말에 50\% 할인했을 경우 연말에 옷의 가격은?}

\[
100 \times (1 + 0.3) \times 0.5 = 100 \times 1.3 \times 0.5 = 65
\]




